\scanpage{7}
\mspage{\handwritten{18}}
\diagram{\[
\begin{tikzcd}[ampersand replacement=\&]
\mathrm{Gr}^{\handwritten{i}}(X_{\bar t})_{\handwritten{\mathbb Q}}
 \handwritten{\arrow[r,"{\mathrm{cl}}"]\arrow[d]} \&
 H^{2i}(X_{\bar t},\underline{Q}_{\unsure{l|\ell}}(i))
 \handwritten{\arrow[d,"\wr"]}\\
\mathrm{Gr}^{\handwritten{i}}(X_{\bar s})_{\handwritten{\mathbb Q}}
 \handwritten{\arrow[r,"{\mathrm{cl}}"]} \&
 H^{2i}(X_{\bar s},\underline{Q}_{l}(i))\;.
\end{tikzcd}\tag*{(7.13.10\deleted{\unsure{9}}\handwritten{)}}
\]}

\deleted{7.13.11.} \uline{Remarque\deleted{s}} 7.13.11. La commutativité des deux derniers diagrammes \deleted{commutatifs}, qu'on peut exprimer en disant que la spécialisation des cycles \unsure{algébriqes|algébriques} est compatible avec la formation des classes de cohomologie $l$-adiques as\unsure{\^{s}|s}ociées, n'est pas très satisfaisante sous la forme énoncée, puisqu'on \deleted{\unsure{ne}} a été obligé de remplacer les coéfficients $\underline{\mathbb Z}_{l}(i)$ par $\underline{Q}_{l}(i)$, c'est à dire de négliger la torsion. Pour avoir un énoncé qui ne néglige pas la torsion, \deleted{on est}\inserted{il semble qu'on sera} obligé de travailler avec l'anneau de Chow, puisque c'est sur ce dernier que ce trouve défini naturellement l'homomorphisme $\mathrm{cl}$ comme homomorphisme à valeurs dans $H^{2*}(-,\underline{Z}_{l}(\unsure{*|i}))$, \deleted{\unsure{Malheureusement,}} comparer (XIV 4.1). Malheureusement, on n'a pas développé pour le moment une théorie de la spécialisation pour l'anneau de Chow\deleted{\unsure{,}} (cf. 7.1\deleted{\unsure{5|3}}\handwritten{4}), qui permettrait d'écrire \deleted{un diagramme}\inserted{des carrés} \deleted{\unsure{a}}commutatifs analogues à (7.13.9) et (7.13.10) en termes des anneaux de Chow et de la cohomologie à coéfficients $\underline{Z}_{l}(*)$. Tout ce qu'on peut dire\inserted{\handwritten{ lorsque $f:X\longrightarrow S$ est lisse et propre\unsure{,}}}, c'est qu'on a un carré commutatif
\diagram{\[
\begin{tikzcd}[ampersand replacement=\&]
Z^i(X_t) \handwritten{\arrow[r,"{\mathrm{cl}}"]\arrow[d,"\sigma"']} \&
 H^{2i}(X_t,\underline{Z}_{l}(i)) \handwritten{\arrow[d,phantom,"{\deleted{\downarrow}}" description,near end]}\\
Z^i(X_s) \handwritten{\arrow[r,"{\mathrm{cl}}"]} \&
 H^{2i}(X_s,\underline{Z}_{l}(i))\handwritten{\arrow[u]}\;,
\end{tikzcd}\tag*{(7.13.12)}
\]}
\deleted{et un carré} analogue à (7.13.9), et un carré correspondant analogue à (7.13.10) qu'on laisse au\deleted{\unsure{tre|x}} lecteur \inserted{\handwritten{le soin}} d'écrire, où $Z^i$ désigne le groupe des cycles algébriques (modulo rien du tout) de codimension $i$, \deleted{et} où les flêches horizontales sont les homomorphismes ``classe de cohomologie associée'' définie dans (SGA 5 IV), et où $\handwritten{\sigma}$ désigne l'homomorphisme de spécialisation des cycles algébriques, associant à un cycle $Z_t$ le cycle $Z_s$, où $Z$ est
% archive sheet number: 6
